\documentclass{scrbook}
\usepackage[utf8]{inputenc}

%General packages for basic document setup
\usepackage{amsmath,amsfonts, amsthm, graphicx}
\usepackage[margin=1.5cm,bottom=3cm, top=3cm]{geometry}

% Random text generator package
\usepackage{lipsum}

% Hyperlink package
\usepackage{hyperref}
\hypersetup{
  colorlinks=true,    %active colors for links
  linkcolor=blue,     %internal link color
  urlcolor=blue,      %external link color
  pdftitle={Template},%title of document (usually for pdf readers to display)
}

% Verbatism package : used for code display 
\usepackage{fancyvrb}
\usepackage{minted} %Interpret code and display colors aswell.

% Header and footer package
\usepackage{fancyhdr, lastpage}
% Header and footer parameters
% Some basic elements are \thepage for the number of the current page,
% \pageref{LastPage} for the number of the last page (add * to remove 
% the hyperlink.
% \rightmark for the current section, \leftmark for the current chapter,
\pagestyle{fancy}


%Chapter-Style : Only use with documentclass that require chaters !
\usepackage[Glenn]{fncychap}
\ChNameVar{\bfseries\Large\sffamily}
\ChTitleVar{\bfseries\Large\rmfamily}

%Colors package
\usepackage{xcolor}


% package for better refs
\usepackage[nameinlink]{cleveref}
% Some parameters to recognize tcbtheorem boxes as theorems, propositions, ...
\crefformat{theoremenv}{#2Théorème~#1#3}
\Crefformat{theoremenv}{#2Théorème~#1#3}

\crefformat{propositionenv}{#2Proposition~#1#3}
\Crefformat{propositionenv}{#2Proposition~#1#3}

\crefformat{lemmaenv}{#2Lemme~#1#3}
\Crefformat{lemmaenv}{#2Lemme~#1#3}

\crefformat{corollaryenv}{#2Corollaire~#1#3}
\Crefformat{corollaryenv}{#2Corollaire~#1#3}

\crefformat{definitionenv}{#2Définition~#1#3}
\Crefformat{definitionenv}{#2Définition~#1#3}

\crefformat{exampleenv}{#2Exemple~#1#3}
\Crefformat{exampleenv}{#2Exemple~#1#3}

\crefformat{remarkenv}{#2Remarque~#1#3}
\Crefformat{remarkenv}{#2Remarque~#1#3}

\newcommand{\creflastconjunction}{ et }

%Figure package
\usepackage{tikz}

%Boxes packages, for theorems, definitions, ...
%REQUIRES tikz in order to work !
\usepackage[most]{tcolorbox}
\tcbuselibrary{breakable}
% if only tcbset is used, it is not possible to make nested environments
% Hence the use of tcbsetforeverylayer
\tcbsetforeverylayer{enhanced jigsaw,breakable,fonttitle=\bfseries}

% ----------- New environments ----------

\usepackage{aliascnt}
% Theorems

\newtheoremstyle{theoremstyle}% name of the style to be used
  {\topsep}% measure of space to leave above the theorem. E.g.: 3pt
  {\topsep}% measure of space to leave below the theorem. E.g.: 3pt
  {\itshape}% name of font to use in the body of the theorem
  {0pt}% measure of space to indent
  {\bfseries}% name of head font
  {}% punctuation between head and body
  { }% space after theorem head; " " = normal interword space
  {
    \textcolor{red}
    {\thmname{#1} \thmnumber{#2}{\thmnote{#3}}}
  }

\theoremstyle{theoremstyle}
\newtheorem{theoremenv}{Théorème}[chapter]
\newenvironment{theorem}[1]{
  \begin{tcolorbox}[frame hidden, boxrule=0pt ,
    borderline west={2pt}{0pt}{red}, colback=red!10!white]{}
    \begin{theoremenv}{#1}
  }{
  \end{theoremenv}
\end{tcolorbox}
}

\newtheoremstyle{proofstyle}% name of the style to be used
  {\topsep}% measure of space to leave above the theorem. E.g.: 3pt
  {\topsep}% measure of space to leave below the theorem. E.g.: 3pt
  {}% name of font to use in the body of the theorem
  {0pt}% measure of space to indent
  {\bfseries}% name of head font
  {\textcolor{red!70!white}:}% punctuation between head and body
  { }% space after theorem head; " " = normal interword space
  {
    \textcolor{red!70!white}
    {\thmname{#1}}
  }

\theoremstyle{proofstyle}
\newtheorem{proofenv}{Preuve}
\renewenvironment{proof}{
  \vspace*{-0.74cm}
  \begin{tcolorbox}[frame hidden, boxrule=0pt ,
    borderline west={2pt}{0pt}{purple!60!white}, colback=red!10!white]{}
  \begin{proofenv}
  }{
  \qed
  \end{proofenv}
  \end{tcolorbox}
}

%Propositions
\newaliascnt{propositionenv}{theoremenv}
\newtheoremstyle{propositionstyle}% name of the style to be used
  {\topsep}% measure of space to leave above the proposition. E.g.: 3pt
  {\topsep}% measure of space to leave below the proposition. E.g.: 3pt
  {\itshape}% name of font to use in the body of the proposition
  {0pt}% measure of space to indent
  {\bfseries}% name of head font
  {}% punctuation between head and body
  { }% space after proposition head; " " = normal interword space
  {
    \textcolor{orange}
    {\thmname{#1} \thmnumber{#2}{\thmnote{#3}}}
  }

\theoremstyle{propositionstyle}
\newtheorem{propositionenv}[propositionenv]{Proposition}
\newenvironment{proposition}[1]{
  \begin{tcolorbox}[frame hidden, boxrule=0pt ,
    borderline west={2pt}{0pt}{orange}, colback=orange!10!white]{}
    \begin{propositionenv}{#1}
  }{
  \end{propositionenv}
\end{tcolorbox}
}
\aliascntresetthe{propositionenv}

%Lemmas
\newaliascnt{lemmaenv}{theoremenv}
\newtheoremstyle{lemmastyle}% name of the style to be used
  {\topsep}% measure of space to leave above the lemma. E.g.: 3pt
  {\topsep}% measure of space to leave below the lemma. E.g.: 3pt
  {\itshape}% name of font to use in the body of the lemma
  {0pt}% measure of space to indent
  {\bfseries}% name of head font
  {}% punctuation between head and body
  { }% space after lemma head; " " = normal interword space
  {
    \textcolor{yellow!60!black}
    {\thmname{#1} \thmnumber{#2}{\thmnote{#3}}}
  }

\theoremstyle{lemmastyle}
\newtheorem{lemmaenv}[lemmaenv]{Lemme}
\newenvironment{lemma}[1]{
  \begin{tcolorbox}[frame hidden, boxrule=0pt ,
    borderline west={2pt}{0pt}{yellow}, colback=yellow!10!white]{}
    \begin{lemmaenv}{#1}
  }{
  \end{lemmaenv}
\end{tcolorbox}
}
\aliascntresetthe{lemmaenv}

%Corollaries
\newaliascnt{corollaryenv}{theoremenv}
\newtheoremstyle{corollarystyle}% name of the style to be used
  {\topsep}% measure of space to leave above the corollary. E.g.: 3pt
  {\topsep}% measure of space to leave below the corollary. E.g.: 3pt
  {\itshape}% name of font to use in the body of the corollary
  {0pt}% measure of space to indent
  {\bfseries}% name of head font
  {}% punctuation between head and body
  { }% space after corollary head; " " = normal interword space
  {
    \textcolor{yellow!60!black}
    {\thmname{#1} \thmnumber{#2}{\thmnote{#3}}}
  }

\theoremstyle{corollarystyle}
\newtheorem{corollaryenv}[corollaryenv]{Corollaire}
\newenvironment{corollary}[1]{
  \begin{tcolorbox}[frame hidden, boxrule=0pt ,
    borderline west={2pt}{0pt}{yellow}, colback=yellow!10!white]{}
    \begin{corollaryenv}{#1}
  }{
  \end{corollaryenv}
\end{tcolorbox}
}
\aliascntresetthe{corollaryenv}

%Definitions
\newaliascnt{definitionenv}{theoremenv}
\newtheoremstyle{definitionstyle}% name of the style to be used
  {\topsep}% measure of space to leave above the definition. E.g.: 3pt
  {\topsep}% measure of space to leave below the definition. E.g.: 3pt
  {}% name of font to use in the body of the definition
  {0pt}% measure of space to indent
  {\bfseries}% name of head font
  {}% punctuation between head and body
  { }% space after definition head; " " = normal interword space
  {
    \textcolor{green!80!black}
    {\thmname{#1} \thmnumber{#2}{\thmnote{#3}}}
  }

\theoremstyle{definitionstyle}
\newtheorem{definitionenv}[definitionenv]{Définition}
\newenvironment{definition}[1]{
  \begin{tcolorbox}[frame hidden, boxrule=0pt ,
    borderline west={2pt}{0pt}{green}, colback=green!10!white]{}
    \begin{definitionenv}{#1}
  }{
  \end{definitionenv}
\end{tcolorbox}
}
\aliascntresetthe{definitionenv}

%Examples
\newaliascnt{exampleenv}{theoremenv}
\newtheoremstyle{examplestyle}% name of the style to be used
  {\topsep}% measure of space to leave above the example. E.g.: 3pt
  {\topsep}% measure of space to leave below the example. E.g.: 3pt
  {}% name of font to use in the body of the example
  {0pt}% measure of space to indent
  {\bfseries}% name of head font
  {}% punctuation between head and body
  { }% space after example head; " " = normal interword space
  {
    \textcolor{blue!80!white}
    {\thmname{#1} \thmnumber{#2}{\thmnote{#3}}}
  }

\theoremstyle{examplestyle}
\newtheorem{exampleenv}[exampleenv]{Exemple}
\newenvironment{example}[1]{
  \begin{tcolorbox}[frame hidden, boxrule=0pt ,
    borderline west={2pt}{0pt}{blue}, colback=blue!10!white]{}
    \begin{exampleenv}{#1}
  }{
  \end{exampleenv}
\end{tcolorbox}
}
\aliascntresetthe{exampleenv}

%Remarks
\newaliascnt{remarkenv}{theoremenv}
\newtheoremstyle{remarkstyle}% name of the style to be used
  {\topsep}% measure of space to leave above the remark. E.g.: 3pt
  {\topsep}% measure of space to leave below the remark. E.g.: 3pt
  {}% name of font to use in the body of the remark
  {0pt}% measure of space to indent
  {\bfseries}% name of head font
  {}% punctuation between head and body
  { }% space after remark head; " " = normal interword space
  {
    \textcolor{violet!80!white}
    {\thmname{#1} \thmnumber{#2}{\thmnote{#3}}}
  }

\theoremstyle{remarkstyle}
\newtheorem{remarkenv}[remarkenv]{Remarque}
\newenvironment{remark}{
  \begin{tcolorbox}[frame hidden, boxrule=0pt ,
    borderline west={2pt}{0pt}{violet}, colback=violet!10!white]{}
    \begin{remarkenv}{}
  }{
  \end{remarkenv}
\end{tcolorbox}
}
\aliascntresetthe{remarkenv}

% --- Exercises and solutions environments ---- 
%Use this package for correct use of exercise and answer environment.
\usepackage{exercise}

% Config of the exercise package
\renewcounter{Exercise}[chapter] %counter for the exercises that resets on each chapter.
\renewcounter{Answer}[chapter]

% Header for the exercises
\def\ExerciseHeader{\noindent\ExerciseHeaderDifficulty\textcolor{brown}{\textbf{Exercice \ExerciseHeaderNB} : \ExerciseHeaderTitle} \\ }

% Header for the answers
\def\AnswerHeader{\noindent\textcolor{brown}{\textbf{Solution de l'exercice \ExerciseHeaderNB}}\\ }

%Exercise environment
\newenvironment{exo}[1]{
  \stepcounter{Exercise}
  \begin{tcolorbox}[frame hidden, boxrule=0pt ,
    borderline west={2pt}{0pt}{brown}, colback=brown!10!white]{}
  \begin{Exercise}[title={#1},label={\the\value{chapter}-\the\value{Exercise}}, number={\the\value{Exercise}},
    counter={Exercise}]
  }{
  % Next line is to create a hyperref to the solution
  % [Solution \refAnswer{\ExerciseLabel}]
  \end{Exercise} 
\end{tcolorbox}
}
%Solution environment
\newenvironment{ans}{
  \begin{tcolorbox}[frame hidden, boxrule=0pt ,
    borderline west={2pt}{0pt}{brown!50!white}, colback=brown!5!white]{}
  \begin{Answer}[ref=\ExerciseLabel]
}{\end{Answer} 
\end{tcolorbox}
}

%%%%% USAGE ADVICE : ALWAYS USE SOLUTION ENVIRONMENT RIGHT AFTER AN EXERCISE ENVIRONMENT OTHERWISE,
%%%%% IT WILL BREAK THE NUMBERING SYSTEM.

% Line spacing package
\usepackage{setspace}
\onehalfspacing

% package used for multicolumns 
\usepackage{multicol}


\sloppy
\usepackage[final, nopatch=footnote]{microtype} % the option is to avoid the 
% Warning : Unable to apply patch 'footnote' 

% ----------- New commands ------------ %


\newcommand*{\blank}{\hspace*{1pt}}
\newcommand*{\MnR}[0]{\mathscr{M}_{n}(\mathbb{R})}
\newcommand*{\GLnR}[0]{\mathbf{GL}_{n}(\mathbb{R})}
\newcommand*{\transpose}[1]{{}^t #1}
\newcommand*{\GL}[0]{\textup{GL}}
\newcommand{\Alt}{\mathfrak{A}}
\newcommand*{\Sym}[0]{\mathfrak{S}}
\newcommand{\W}[0]{\textup{W}}
\newcommand{\Stab}[0]{\textup{Stab}}
\newcommand{\Fix}[0]{\textup{Fix}}
\newcommand*{\rg}[0]{\textup{rg}}
\newcommand*{\Im}[0]{\textup{Im}}
\newcommand*{\Ker}[0]{\textup{Ker}}
\newcommand*{\Coker}[0]{\textup{Coker}}
\newcommand*{\Tr}[0]{\textup{Tr}}
\newcommand*{\sinc}[0]{\textup{sinc}}
\newcommand*{\abs}[1]{\left| #1 \right|}
\newcommand*{\norm}[1]{\left\lVert #1 \right\rVert}
\newcommand*{\indicatrice}[0]{\mathds{1}}
\newcommand*{\Vect}[1]{\textup{Vect}\left\{#1\right\}}
 %All needed information are in this document

%Global information about the document
\def\THETITLE{Notes de cours 0}
\def\THEAUTHOR{Romain Delaunay}

\title{\THETITLE}
\author{\THEAUTHOR}
\date{}


% Header
\lhead{\THETITLE}
\rhead{\rightmark}

% Foot
\lfoot{\THEAUTHOR}
\cfoot{\thepage \ / \pageref*{LastPage} }
\rfoot{}

\hypersetup{
  pdftitle={NotesDeCours0}
}
%------------- Document --------------%

\begin{document}

\addto\extrasfrench{\protected\edef:{\unexpanded\expandafter{:}}}
\selectlanguage{french}

\maketitle

\tableofcontents

\chapter{Test de chapitre}

Ceci est un exemple en français.

\section{Quotient}

\begin{definition}{(Quotient)}
    On appelle quotient de l'ensemble $X$ toute paire $(Y,\pi )$ avec $Y$ un ensemble quelconque, et $\pi :X \to Y$ 
    une application surjective. On dit que ce quotient est associé à la relation $\sim$ si : 
    $\forall x,x' \in X, x \sim x' \Leftrightarrow \pi (x) = \pi (x')$. 
\end{definition}

Dans la suite, si le contexte le permet, on pourra noter simplement $Y$ le quotient.

\begin{example}
  \blank
  \begin{enumerate}
    \item Pour tout ensemble $X$ et toute relation d'équivalence $\sim$ sur $X$, on note $\faktor{X}{\sim} := \{ 
      \textup{cl}(x) ~|~ x \in X \}$ où $\textup{cl}(x) := \{ y \in X ~|~ x \sim y \}$ est la classe 
      d'équivalence de $x$. Alors $(\faktor{X}{\sim}, \textup{cl})$ est un quotient de $X$ associé à la relation $\sim$.
    \item On pose sur $\mathbb{N}^{2}$ la relation d'équivalence définie par : $\forall (a,b), (x,y) \in \mathbb{N}^{2}, 
      (a,b) \sim (x,y) \Leftrightarrow a+y = x + b$. On note alors le quotient $\mathbb{Z}$. Évidemment, cette notation
      n'est pas choisie au hasard puisque le but ici est de construire les entiers relatifs partant des entiers naturels.
      Il est laisser (pour le moment) aux lectrices et lecteurs le soin de trouver une représentation graphique de 
      ce quotient et d'en déduire le lien avec la structure naturelle et usuelle de $\mathbb{Z}$.
  \end{enumerate}
\end{example}   

\begin{theorem}{(Propriété universelle des quotients)}
    \label{thm:ppuquotient}
    Soit $\pi : X \to Y$ un quotient associé à $\sim$ et soit $f :  X \to Z$ telle que : 
    \[
      \forall x,x', x \sim x' \Rightarrow f(x)=f(x')
    \]
    Alors : 
    \[
      \exists ! \varphi : Y \to Z, f = \varphi \circ \pi
    \]

    Autrement dit, il existe un unique $\varphi : Y \to Z$ telle que le diagramme suivant commute : 
    \begin{center}
      \begin{tikzcd}[column sep=2cm, row sep=2cm]
        X \arrow[r,"f"] \arrow[d,"\pi"] & Y \\
        Z \arrow[ur,dashed,dash,yshift=-3pt,below,"\varphi"'] \arrow[ur]
      \end{tikzcd}
    \end{center}

    De plus, 

    \begin{enumerate}
      \item $\varphi$ est surjective si, et seulement si, $f$ l'est.
      \item $\varphi$ est injective si, et seulement si : $\forall x,x', x \sim x' \Leftrightarrow f(x) =f(x')$.
    \end{enumerate}

    Réciproquement, si $\varphi$ existe, alors $\forall x,x', x \sim x' \Rightarrow f(x) = f(x')$.
\end{theorem}

\begin{proof}
  L'idée est que, puisque tout élément d'une classe d'équivalence est envoyé sur la même image, 
  prendre un représentant pour chaque classe et lui associer son image constituera le morphisme
  $\varphi$ recherché. 
  Par l'axiome du choix, on peut se donner $\sigma : Y \to X$ telle que $\pi \circ \sigma = \id_Y$.

  On pose alors $\varphi := f \circ \sigma$ (cette construction dit bien qu'à chaque classe 
  d'équivalence, on associe l'image d'un de ses éléments). 

  On a $\varphi \circ \pi = f \circ \sigma \circ \pi$. 
  Or, pour $x \in X$, on a $\sigma \circ\pi (x) \sim x$. 
  En effet, $\pi \left( \sigma \circ\pi (x) \right) = \left(\pi \circ \sigma\right) \circ\pi(x) =\pi (x)$. 
  Ainsi, par l'hypothèse donnée sur $f$, $f( \sigma \circ\pi (x)) = f(x)$, et ce, pour tout $x \in X$.

  On a bien montré que $f = \varphi \circ\pi$.
  Ceci montre l'existence d'un tel $\varphi$. Il reste alors à montrer l'unicité.

\end{proof}

\begin{proofr}
  \underline{(Unicité)} : 
    Soit $\tilde{\varphi} : Y \to Z$ telle que $f = \tilde{\varphi} \circ \pi $. 
    Montrons que $\varphi = \tilde{\varphi}$.
    
    Soit $y \in Y$, par surjectivité de $\pi $, on peut se donner $x \in X$ tel que $y = \pi (x)$.

    On obtient donc que : 
    \begin{align*}
      \varphi(y)  &= \varphi(\pi (x)) \\
                  &= f(x) \\
                  &= \tilde{\varphi}(\pi (x)) \\
                  &= \tilde{\varphi}(y) 
    \end{align*}
    et ce, pour tout $y \in Y$. Ceci montre bien que $\tilde{\varphi} = \varphi$.
        D'où l'unicité.

\end{proofr}

\begin{proofe}
    
\end{proofe}
  \begin{proofe}
  Montrons désormais le résultat portant sur la surjectivité de $\varphi$.
  \end{proofe}
  
  \begin{proofr}
      \underline{(Surjectivité)} : On raisonne par double implication.
  \end{proofr}

  \begin{proofrr}
      \underline{($\Rightarrow$)} : Si $\varphi$ est surjective, $f$ l'est également comme composée de telles fonctions.
            D'où l'implication directe. 

        \underline{($\Leftarrow$)} : Supposons que $f$ soit surjective. Pour $z \in Z$, on se donne $x \in X$ tel que $f(x)=z$.
        Alors $\varphi(\pi (x)) = f(x) = z$. D'où la surjectivité de $\varphi$.
            D'où l'implication réciproque.
  \end{proofrr}
    \begin{proofr}
      D'où l'équivalence portant sur la surjectivité de $\varphi$.
    \end{proofr}


\begin{proofe}
  Montrons désormais l'équivalence portant sur l'injectivité de $\varphi$.
\end{proofe}

  \begin{proofr}
    \underline{(Injectivité)} : 
      Une fois n'est pas coutume, nous allons raisonner par équivalence en utilisant
      notamment le fait que $\pi $ soit surjective.

      \[
        \begin{WithArrows}
          \varphi \text{ est injective} &\Leftrightarrow \forall y,y' \in Y, \varphi(y) =
                                            \varphi(y') \Rightarrow y=y' \curvearrow{$\pi $ est surj.} \\
                                        & \Leftrightarrow \forall x,x' \in X, \varphi(\pi (x)) =
                                            \varphi(\pi (x')) \Rightarrow \pi (x) = \pi (x') \\
                                        & \Leftrightarrow \forall x,x' \in X, f (x) = f(x') 
                                            \Rightarrow x \sim x' \curvearrow{autre imp. déjà donnée} \\
                                        & \Leftrightarrow \forall x,x' \in X, f (x) = f(x')   
                                            \Leftrightarrow x \sim x' 
        \end{WithArrows}
      \]

          D'où l'équivalence portant sur l'injectivité de $\varphi$.
  \end{proofr}

\begin{proofe}

  Enfin, si réciproquement une telle fonction $\varphi$ existe, alors pour $x,x' \in X$ tels que 
  $x \sim x'$, $f(x) = \varphi(\pi (x)) = \varphi(\pi (x')) = f(x')$.

\qed

\end{proofe}

\begin{remark}
    La partie unicité peut se faire plus rapidement en invoquant l'inversibilité à droite de $\pi $ par sa surjectivité.
    De plus, il est possible de se passer de l'axiome du choix en montrant que 
    $\{ (y,z) ~|~ \forall x \in X , \pi (x) = y \Rightarrow f(x)=z\}$ est un graphe.
\end{remark}

\begin{example}
  \label{ex:uniquequotient}
\blank
  \begin{enumerate}
    \item 
      Pour tout quotient $(Y,\pi )$ de $X$ associé à $\sim$, $\pi $ passe au quotient par $(\faktor{X}{\sim},\textup{cl}$).
    \item Soit $X$ un ensemble et $d : X \times X \to \mathbb{R}_+$ une application vérifiant :

    \begin{enumerate}
      \item $\forall x,y \in X, d(x,y) = d(y,x)$
      \item $\forall x,y,z \in X, d(x,z) \leq d(x,y) + d(y,z)$.
      \item $\forall x,y \in X, d(x,y) = 0 \Leftrightarrow x \sim y$
    \end{enumerate}

    Alors $d$ induit une application quotient $\tilde{d} : \faktor{X}{\sim }  \times \faktor{X}{\sim } \to \mathbb{R}_+$
    qui définit alors une distance sur $\faktor{X}{\sim }$.
  \end{enumerate}
\end{example}

L' \cref{ex:uniquequotient} de cet exemple montre que si un quotient est associé à une certaine
relation d'équivalence $\sim$, alors on peut factoriser (ou passer au quotient) par $\faktor{X}{\sim}$.
C'est en réalité le seul quotient (à bijection près) associé à $\sim$ comme le montre le corollaire suivant.

\begin{corollary}{}
  \label{cor:uniquequotient}
  Soient $\pi _1 :X \to Y_1$, $\pi _2 : X \to Y_2$ deux quotients associés à la relation d'équivalence $\sim$.
  Alors il existe une unique bijection $\psi : Y_1 \to Y_2$ telle que $\pi _2 = \psi \circ \pi _1$.
\end{corollary}
\begin{proof}
    On regarde le diagramme commutatif suivant : 

\begin{center}
    
  \begin{tikzcd}[column sep=1.5cm, row sep=1.5cm]
        & Y_1 \arrow[dr, "\exists ! \id _{Y_1}"] & \\
    Y_2 \arrow[ur,"\exists ! \varphi"] \arrow[ur,dashed,dash,yshift=-3pt] & X \arrow[u," \pi_1 "] \arrow[r,"\pi _1"] 
                        \arrow[d,"\pi _2"] \arrow[l,"\pi _2"] & Y_1 \arrow[dl,dashed,dash,yshift=-3pt] \arrow[dl, "\exists ! \psi "]\\
                                       & Y_2 \arrow[ul,"\exists ! \id _{Y_2}"] 
  \end{tikzcd}  

\end{center}

En effet, en prenant le quotient $(Y_1,\pi _1)$ de $X$ auquel on applique la propriété universelle des
quotients à $\pi _2$, on obtient le triangle inférieur droite. De même pour $(Y_2,\pi _2)$ avec $\pi _1$,
on obtient le triangle supérieur gauche. Ceci donne l'existence et l'unicité des applications $\psi $ et $\varphi$.
On note enfin qu'elles sont bien inverses l'une de l'autre via la propriété universelle appliquée à
$(Y_1,\pi _1)$ sur $\pi _1$ cette fois, et en utilisant le fait que $\id_{Y_1}$ convient (comme application unique
par l'unicité procurée par la propriété universelle des quotients encore une fois) pour l'application quotient de $\pi _1$,
et de même pour $(Y_2, \pi _2)$ appliqué à $\pi _2$.
Dès lors, par unicité, on obtient que : $\psi \circ \varphi = \id_{Y_1}$ et $\varphi \circ \psi = \id_{Y_2}$.

\qed
\end{proof}

\begin{example}
    Soient $\pi _1 : X_1 \to Y_1$ et $\pi _2 : X_2 \to Y_2$ deux quotients associés respectivement
    à $\sim _1$ et $\sim _2$.
    On définit sur $X_1 \times X_2$ la relation d'équivalance produit $\sim$ par :

    \[
      \forall m_1,m_1' \in X_1 ,m_2,m_2' \in X_2, (m_1,m_2) \sim (m_1',m_2') \Leftrightarrow m_1 \sim_1 m_1' \wedge m_2 \sim_2 m_2'
    \]  

    et $\pi : X_1 \times X_2 \to \faktor{(X_1 \times X_2)}{\sim }$ la surjection canoniquement associée.

    Puisque 
    \begin{align*}
      \forall m_1,m_1' \in X_1,~ m_2, m_2' \in X_2,~ \pi (m_1,m_2) = \pi (m_1',m_2') \Leftrightarrow (m_1,m_2) \sim (m_1',m_2') \\ 
      \Leftrightarrow \begin{cases}
          m_1 \sim _1 m_1' \\
          m_2 \sim _2 m_2' \\
      \end{cases} \Leftrightarrow 
      \begin{cases}
          \pi _1 (m_1) = \pi _1(m_1') \\ 
          \pi _2 (m_2) = \pi _2(m_2')
      \end{cases}
    \end{align*}

    On observe que $Y_1 \times Y_2$ et $\faktor{X_1 \times X_2}{\sim }$ sont deux quotients associés
    à $\sim $. Ils sont donc en bijection par ce qui précède.

    Cet exemple montre d'une manière générale que, avec les bonnes notations : 

    \[
      \faktor{X_1}{\sim _1} \times \faktor{X_2}{\sim _2} \simeq \faktor{(X_1 \times X_2)}{(\sim _1 \times \sim _2)}
    \]
    
\end{example}

Une autre conséquence directe et fondamentale du \cref{thm:ppuquotient} est le double passage au quotient, 'au départ et à l'arrivée', via le 

\begin{corollary}{}
  \label{cor:doublequotient}
  Soient $\pi _{1}:X_{1} \to Y_{1}, \pi _2 : X_2 \to Y_2$ deux quotients associés respectivement à
  $\sim_1$ et $\sim_2$, ainsi que $f : X_1 \to X_2$, une fonction vérifiant : 
  $\forall x,x' \in X_1 , x \sim_1 x' \Rightarrow f(x) \sim_2 f(x')$.
  Alors : $\exists ! \tilde{f} : Y_1 \to Y_2 , \pi _2 \circ f = \tilde{f} \circ \pi _1$.
  Cette condition se traduit par la commutativité du diagramme suivant : 

  \begin{center}

    \begin{tikzcd}[column sep=2cm, row sep=2cm]
      X_1 \arrow[r,"f"] \arrow[d,"\pi_1"] & X_2 \arrow[d,"\pi_2"] \\
      Y_1  \arrow[r,dashed,dash,yshift=-4pt,"\exists ! \tilde{f}"'] \arrow[r] & Y_2
    \end{tikzcd}
      
  \end{center}
\end{corollary}

\begin{proof}
  Il suffit d'appliquer le \cref{thm:ppuquotient} à $\pi _2 \circ f$.

  \qed
\end{proof}

On notera qu'il est également possible de donner une condition nécessaire et suffisante pour la 
surjectivité de $\tilde{f}$, qui ne demande pas nécessairement que $f$ le soit.

\begin{example}
  \label{ex:Zmonoid}
  L'application $+ : \mathbb{N}^2 \to \mathbb{N}$ descend en une application $\bar{+} : \mathbb{Z}^2 \to \mathbb{Z}$
  en identifiant le quotient produit et le produit de quotient.
  On pourra vérifier que $\bar{+}$ est associative, commutative, de neutre $\pi ((0,0))$,
  et telle que pour tout élément $\pi ((a,b)) \in \mathbb{Z}$, $\pi ((a,b)) + \pi ((b,a)) = \pi ((a+b,a+b)) = \pi ((0,0))$.
  Enfin, on peut remarquer que $\mathbb{N}$ s'injecte naturellement dans $\mathbb{Z}$ via $\iota : n \mapsto \pi ((n,0))$.
  On obtient de telle manière un groupe abélien, partant d'un monoïde commutatif.
\end{example}

Cette section est limitée dans les exemples par le manque de structure des objets que nous étudions.
Quand bien même la notion de passage au quotient est fondamentale en algèbre, il est plus que nécessaire
de poursuivre cette étude en observant les liens avec des objets ayant plus de structure.

\nocite{*}
\bibliography{bibli}{}
\bibliographystyle{plain}

\end{document}



